\chapter{Atoms, bonds, periodic table}
\label{vol05:ch01}

\epigraph{The properties of the elements are in periodic dependence
upon their atomic weights.}{Dmitri Mendeleev, presenting the periodic
system to the Russian Chemical Society, March 1869, in the German
abstract \cite[pp.~405--406]{hist:mendeleev1869}.}

\chapmeta{Half-life: foundational. Archetypes: scaling (atomic to macroscopic), interface. Exercise target: 28. Project track: household identification. Hazard class: low (visual and tactile inspection of common materials). Reader path default: standard, with sections explicitly tagged. Named cases: none required in this chapter; the bond-level failure modes are illustrated by mechanism rather than by named accident.}

\noindent
Every macroscopic property of a piece of matter is set by the way its
atoms bond, the geometry in which those bonds arrange them, and the
defects that interrupt the geometry. The chapter develops the chain
from electron configuration through the periodic table to the four
canonical bond types, then to the crystalline and amorphous
arrangements that bonds produce, and closes by reading mechanical and
electrical properties off the bond budget. Later chapters in the volume
specialise the chain to particular materials classes; the present
chapter establishes the vocabulary that those later chapters reuse.

\section{Atomic structure}\pathtag{core}

An atom is a small, dense nucleus of protons and neutrons surrounded
by a cloud of electrons. The nucleus carries essentially all the mass
(the proton and neutron each weigh roughly $1.67 \times 10^{-27}\,
\text{kg}$, the electron $9.11 \times 10^{-31}\,\text{kg}$, a factor of
$1836$ smaller). The nucleus carries all the positive charge ($+Ze$
for an atom of atomic number $Z$, with $e = 1.602 \times 10^{-19}\,
\text{C}$). The electrons carry the negative charge that balances it,
and they occupy a volume tens of thousands of times larger than the
nucleus itself: the nuclear radius is of order $10^{-15}\,\text{m}$,
the atomic radius of order $10^{-10}\,\text{m}$. Almost everything an
engineer measures about a piece of matter at the laboratory bench is a
measurement of the electron cloud; the nucleus determines the
chemistry by setting $Z$ and the gravitational and inertial properties
by setting the mass.

The electron cloud is not a classical orbit. It is the probability
distribution of a quantum-mechanical wave that solves the
Schr\"odinger equation in the potential of the nucleus. For the simple
case of one electron around a point-charge nucleus (the hydrogen atom),
the solutions are the familiar hydrogenic orbitals, indexed by four
quantum numbers: the principal quantum number $n = 1, 2, 3, \ldots$,
which sets the shell; the orbital angular-momentum quantum number
$\ell = 0, 1, \ldots, n-1$, which sets the subshell and is labelled
$s, p, d, f$ in the conventional spectroscopic notation; the magnetic
quantum number $m_\ell = -\ell, \ldots, +\ell$, which labels the
distinct orbitals within a subshell; and the spin quantum number
$m_s = \pm 1/2$. The Pauli exclusion principle says that no two
electrons in the same atom may share the same set of four quantum
numbers. The hydrogenic energy depends only on $n$ for a single
electron, with $E_n = -13.6\,\text{eV}/n^2$, but the energy depends on
both $n$ and $\ell$ in multi-electron atoms because the inner electrons
screen the nuclear charge seen by the outer ones. The engineering
content of the quantum theory is the orbital structure; the
spectroscopic energies of isolated atoms matter only for the few
applications that depend on them directly (atomic clocks, optical
emission spectroscopy, semiconductor band structure).

\engterm{Filling rules}. Multi-electron atoms in their ground state
fill orbitals from low energy to high, in the order set by the
Madelung rule ($n + \ell$ first, then increasing $n$ within each
$n+\ell$ group): $1s$, $2s$, $2p$, $3s$, $3p$, $4s$, $3d$, $4p$, $5s$,
$4d$, $5p$, $6s$, $4f$, $5d$, $6p$, $7s$, $5f$, $6d$, $7p$. The
exceptions are few and known: chromium ($\text{Ar}\,3d^5 4s^1$ rather
than $3d^4 4s^2$), copper ($3d^{10} 4s^1$ rather than $3d^9 4s^2$),
and a handful of heavier elements where the closeness of $d$ and $f$
subshells gives the half- or fully-filled configuration a small
energetic advantage \cite[ch.~9]{text:atkins-physical-chemistry}. The
exceptions matter for the electronic and magnetic properties of those
specific elements (the unusual paramagnetism of copper-rich alloys,
the colour of chromium-containing pigments) and not for the broad
sweep of the periodic system.

\engterm{Valence}. The chemistry of an element is set, to a good
approximation, by its outermost partially filled shell. The
\engterm{valence electrons} are the electrons in that shell. The
periodic table groups elements with the same number of valence
electrons into columns, and the columns predict the dominant ionic
charges, the typical coordination numbers in compounds, and the
qualitative reactivity patterns. The hydrogen atom has one valence
electron ($1s^1$) and forms one bond; carbon has four ($2s^2 2p^2$,
which hybridises to four equivalent $sp^3$ orbitals in saturated
organic compounds) and forms four; the noble gases have full outer
shells ($1s^2$ for helium, $ns^2 np^6$ for the others) and form
essentially no bonds at ordinary conditions.

\section{The periodic table read engineering-first}\pathtag{core}

The periodic table is the engineering reference card for the elements.
The reader who can read the table at a glance can predict, from group
membership and period number, whether an element is metallic, the
ionic charge it favours, the rough magnitude of its atomic radius and
electronegativity, and the qualitative reactivity it shows toward
oxygen, water, and acids. The chapter assumes the reader knows the
table exists and has read the conventional labelling; the next several
paragraphs treat the table as a working chart whose patterns set up
the bonding theory.

\engterm{Periods and groups}. A \engterm{period} is a horizontal row,
labelled $1$ through $7$. The period number is the principal quantum
number $n$ of the highest-occupied shell. A \engterm{group} is a
vertical column, labelled by the IUPAC numbering ($1$ through $18$).
Group $1$ is the alkali metals (one $s$ valence electron, ionises to
$+1$); group $2$ the alkaline earths (two $s$ valence electrons,
ionises to $+2$); groups $3$ through $12$ the transition metals
(partially filled $d$ subshell, variable oxidation states); groups
$13$ through $17$ the $p$-block elements; group $18$ the noble gases.
The block structure ($s$-block, $p$-block, $d$-block, $f$-block) maps
the orbital being filled across the period.

\engterm{Atomic radius and ionic radius}. The atomic radius generally
decreases left-to-right across a period (the added electrons enter the
same shell, while the added nuclear charge pulls all electrons in) and
increases top-to-bottom down a group (each new period adds a shell).
The empirical values cluster: alkali metal atoms have radii in the
$1.5$ to $2.6\,\si{\angstrom}$ range, transition metal atoms in the
$1.3$ to $1.6\,\si{\angstrom}$ range, halogen atoms in the $0.7$ to
$1.4\,\si{\angstrom}$ range. Ionic radii follow the same trends but
shift: cations are smaller than their parent atoms (they have lost
outer electrons), anions are larger (the added electrons increase
inter-electron repulsion and expand the cloud).

\engterm{Ionisation energy}. The energy required to remove the
outermost electron from a neutral atom in the gas phase. The first
ionisation energy of the alkali metals is low ($3.9\,\text{eV}$ for
caesium, $5.4\,\text{eV}$ for lithium) because the single valence
electron is loosely held; for the noble gases it is high ($24.6\,
\text{eV}$ for helium, $15.8\,\text{eV}$ for argon) because the closed
shell is stable. The pattern across a period is increasing, with local
dips at the boron-like and oxygen-like configurations (the half-filled
$p^3$ configuration of group $15$ is unusually stable, and group $16$
loses its first electron more easily than group $15$). The engineer
who reads off ionisation energies from a chart can predict the
qualitative redox behaviour without further calculation.

\engterm{Electron affinity}. The energy released when a neutral atom
gains an electron in the gas phase. The halogens have the largest
electron affinities (chlorine $3.6\,\text{eV}$, fluorine $3.4\,
\text{eV}$); the alkaline earths and noble gases have essentially
zero. The difference between ionisation energy and electron affinity
predicts which atom of an AB pair will be the cation and which the
anion.

\engterm{Electronegativity}. The empirical scale, due to Pauling
\cite[ch.~3]{hist:pauling1960}, of an atom's tendency to attract
electrons in a chemical bond. Pauling's scale assigns fluorine $4.0$
(the most electronegative element), oxygen $3.4$, nitrogen $3.0$,
carbon $2.5$, hydrogen $2.1$, the alkali metals $0.7$ to $1.0$. The
difference in electronegativity between the two atoms in a bond
predicts the bond's ionic character: a difference greater than
approximately $1.7$ produces a predominantly ionic bond (NaCl, with
$\Delta\chi = 2.2$), a difference between $0.5$ and $1.7$ produces a
polar covalent bond (the OH bond in water has $\Delta\chi = 1.3$), a
difference smaller than $0.5$ produces an essentially covalent bond
(the CC bond, $\Delta\chi = 0$, and the CH bond, $\Delta\chi = 0.4$).
The scale is approximate and the boundaries are not sharp, but for
engineering identification of bond type it is the single most useful
empirical chart.

\engterm{The metal-nonmetal line}. A staircase drawn diagonally from
boron to astatine on the right side of the table separates metals
(below and to the left) from non-metals (above and to the right). The
elements adjacent to the line on the metal side (silicon, germanium,
arsenic, antimony, tellurium) are the \engterm{metalloids}; their
bonding is covalent but their electronic structure is on the border
of metallic, and they are the semiconductors of practical engineering
importance.

The periodic table is older than the quantum theory it now organises.
Mendeleev's 1869 announcement \cite{hist:mendeleev1869} arranged the
known elements by atomic weight, identified gaps in the sequence, and
predicted the properties of three then-undiscovered elements (gallium,
germanium, scandium) before they were isolated. The predictions matched
the measured atomic weights, densities, and reactivities within tens
of percent. The success was empirical: the modern quantum
understanding of why the elements are periodic (the filling of $s$,
$p$, $d$, $f$ subshells) came forty years later, with the
Bohr-Sommerfeld atomic model and then the Schr\"odinger equation. The
working engineer's relationship to the table is empirical first and
quantum second: the patterns are robust regardless of which derivation
one uses to defend them.

\begin{archetype}[Scaling: from atomic to macroscopic]
The macroscopic properties of a piece of matter (its density,
stiffness, electrical conductivity, thermal expansion, optical
opacity) arise from properties of single atoms (their mass, their
bond energies, their bond lengths) by aggregation over Avogadro's
number $N_A = 6.022 \times 10^{23}$ atoms per mole. The aggregation is
usually direct, with multiplicative scaling: a mole of carbon has mass
$12\,\text{g}$, and a kilogram of diamond contains $N_A \times 1000/12
\approx 5 \times 10^{25}$ carbon atoms. The bond energy of a single
C-C bond ($\sim 3.6\,\text{eV} = 5.8 \times 10^{-19}\,\text{J}$) times
the bond density in diamond ($\sim 4 \times 10^{29}\,\text{m}^{-3}$,
since each atom has four bonds shared between two atoms) gives a
volumetric energy density of $\sim 1.2 \times 10^{11}\,\text{J}/
\text{m}^3$, which is the order of magnitude of the energy required to
mechanically pulverise diamond at the bond level. The reader who can
move between the atomic and the macroscopic by this kind of bookkeeping
owns the working core of Volume V.
\end{archetype}

\section{Bonding: ionic, covalent, metallic, secondary}\pathtag{core}

Atoms bond when bonding lowers the system's total energy. The four
canonical bond types (ionic, covalent, metallic, secondary) differ in
how the electrons are arranged between the bonded atoms; they share
the same energetic logic. Real bonds in real materials are usually
mixtures of the four types, weighted by the electronegativity
difference and the geometry. The four categories are an engineering
classification of bond character, useful for predicting properties,
not a partition of bonds into disjoint classes.

\subsection{Ionic bonding}

An ionic bond transfers an electron from a low-ionisation-energy atom
to a high-electron-affinity atom and binds the resulting ions by
Coulomb attraction. The classic case is sodium chloride: sodium
($\text{Ne}\,3s^1$, ionisation energy $5.1\,\text{eV}$) loses its
$3s$ electron, chlorine ($\text{Ne}\,3s^2 3p^5$, electron affinity
$3.6\,\text{eV}$) gains it, and the resulting Na$^+$ and Cl$^-$ ions
arrange themselves in a three-dimensional lattice in which each ion is
surrounded by six oppositely charged neighbours. The Coulomb energy of
a pair of point charges $+e$ and $-e$ separated by a distance $r$ is
$-e^2/(4\pi\varepsilon_0 r)$; in a lattice, the energy per ion is the
sum over all other ions, computed by the \engterm{Madelung constant}
$M$ for the lattice geometry. For NaCl the Madelung constant is
$M = 1.7476$, and the lattice energy per ion pair is
\[
U_{\text{ion}} = -\frac{M e^2}{4\pi\varepsilon_0 r_0}\left(1 -
\frac{1}{n}\right),
\]
where $r_0$ is the nearest-neighbour distance and $n$ is the Born
repulsion exponent (typically $5$ to $12$) that accounts for the
short-range electron-cloud overlap repulsion. For NaCl, $r_0 = 2.82\,
\si{\angstrom}$ gives $U_{\text{ion}} \approx -7.9\,\text{eV}/\text{pair}
= -762\,\text{kJ}/\text{mol}$, within a few percent of the measured
lattice energy. The agreement is the canonical demonstration that
ionic bonding is correctly described by Coulomb energy plus a
short-range repulsion.

Ionic bonds are strong (lattice energies of order $500$ to $4000\,
\text{kJ}/\text{mol}$), isotropic (the Coulomb force is central), and
non-directional (an ion bonds to all its neighbours of opposite
charge equally). The resulting solids are hard but brittle: a shear
displacement that brings like ions next to each other produces strong
electrostatic repulsion and the lattice cleaves. Ionic solids
typically have high melting points (NaCl $801\,\degC$, MgO $2852\,
\degC$, Al$_2$O$_3$ $2072\,\degC$), low electrical conductivity in
the solid state (electrons are bound to specific ions, and ions are
not mobile through the lattice at room temperature), and high
electrical conductivity in the molten state or in solution (ions
become mobile).

\subsection{Covalent bonding}

A covalent bond shares a pair of electrons between two atoms. The
shared pair occupies a molecular orbital constructed by combining the
atomic orbitals of the two atoms; in the simplest description the
combination is a linear combination of atomic orbitals (LCAO),
producing a bonding orbital with lower energy than the parent atomic
orbitals and an antibonding orbital with higher energy. Two electrons
fill the bonding orbital and the bond is stable. The hydrogen molecule
H$_2$ has bond energy $4.5\,\text{eV} = 436\,\text{kJ}/\text{mol}$ and
bond length $0.74\,\si{\angstrom}$; the C-C single bond has $3.6\,
\text{eV} = 348\,\text{kJ}/\text{mol}$ and $1.54\,\si{\angstrom}$; the
C-H bond, $4.3\,\text{eV} = 414\,\text{kJ}/\text{mol}$ and $1.09\,
\si{\angstrom}$. Bond energies cluster between $2$ and $10\,\text{eV}$
for single bonds; double and triple bonds add roughly $50\,\%$ and
$100\,\%$ to the energy of the equivalent single bond.

Covalent bonds are directional: the bonding orbital has a specific
spatial orientation set by the parent atomic orbitals. In carbon the
four $sp^3$ hybrid orbitals point to the corners of a tetrahedron, and
diamond accordingly has tetrahedral coordination with a bond angle of
$109.5^{\circ}$. In graphite, three of carbon's four valence electrons
form $sp^2$ hybrid orbitals in a plane (bond angle $120^{\circ}$) and
the fourth $p$ electron forms a delocalised $\pi$ bond above and below
the plane; the result is the layered graphite structure with strong
in-plane covalent bonding and weak inter-plane van der Waals bonding.
The contrast between diamond (the hardest natural material, electrical
insulator, melting point near $4000\,\degC$) and graphite (soft along
the layers, good in-plane electrical conductor, sublimes near
$3650\,\degC$ at one atmosphere) is the canonical demonstration that
the same atoms with the same bond strength can produce wildly
different macroscopic properties through different bond geometry.

The covalent bond is the bond of organic chemistry, of the molecular
solids (ice, sugar, sulphur), of the network solids (diamond, silicon
carbide, boron nitride, fused silica), and of the polymers (Chapter~6
of this volume). The directionality of the bond produces the rich
geometric variety of organic and macromolecular structures; the
strength of the bond is responsible for the structural performance of
the engineering polymers and ceramics.

\subsection{Metallic bonding}

A metallic bond distributes the valence electrons of all atoms over
the entire crystal, producing a ``sea'' of mobile electrons in which
the positively charged ion cores sit. The free-electron model treats
the conduction electrons as a degenerate Fermi gas, with the
Fermi energy $E_F$ set by the electron density:
\[
E_F = \frac{\hbar^2}{2 m_e} (3\pi^2 n)^{2/3},
\]
where $n$ is the conduction-electron number density. For copper
($n = 8.45 \times 10^{28}\,\text{m}^{-3}$, one $4s$ electron per atom),
$E_F \approx 7.0\,\text{eV}$; for aluminium ($n = 1.81 \times 10^{29}\,
\text{m}^{-3}$, three valence electrons per atom), $E_F \approx 11.7\,
\text{eV}$. The bond energy per atom in a metal is typically
comparable to the free-atom ionisation energy (one to several
electron-volts).

Metallic bonds are non-directional (the electron sea spreads
uniformly), strong (cohesive energies of $1$ to $9\,\text{eV}/\text{atom}$),
and tolerant of geometric rearrangement: an atom can be displaced
relative to its neighbours without breaking specific localised bonds,
because the bonding is collective. The resulting solids are
mechanically ductile (they yield by dislocation motion rather than by
cleavage, Chapter~5 of this volume), electrically conductive (the
free electrons carry current with mobilities of $10$ to $100\,
\text{cm}^2/(\text{V}\cdot\text{s})$ at room temperature), thermally
conductive (the same electrons carry heat, with the
Wiedemann-Franz law $\kappa/(\sigma T) \approx 2.45 \times 10^{-8}\,
\text{W}\Omega/\si{\kelvin}^2$ relating the thermal and electrical
conductivities), and optically reflective at visible wavelengths (the
plasma frequency of the electron sea is in the ultraviolet for most
metals, so visible light is reflected). Metallic bonding is the
working bond of structural engineering, of electrical infrastructure,
and of every metal-bodied machine in the engineered world.

\subsection{Secondary bonding}

Secondary bonds are weaker, non-directional bonds that hold molecules
to one another even when the molecules themselves are held together
by covalent bonds. The three principal types:

\engterm{Van der Waals bonds}. Instantaneous dipole-induced-dipole
interactions, arising because the electron clouds of two molecules
fluctuate, and a fluctuation in one induces a complementary fluctuation
in the other. Van der Waals bonds have energies of $0.05$ to $0.4\,
\text{eV}$ per pair and act over distances of $3$ to $5\,
\si{\angstrom}$. They are the bonds that hold molecular solids (argon
ice at $84\,\si{\kelvin}$, solid CO$_2$ at $-78\,\degC$) together; they
are responsible for the inter-layer bonding in graphite that allows the
layers to slide; they govern the boiling points of non-polar liquids.

\engterm{Hydrogen bonds}. A hydrogen atom covalently bonded to a small,
electronegative atom (oxygen, nitrogen, fluorine) carries a
substantial positive charge (the electron of the H-O bond is
displaced toward the oxygen) and forms an attractive interaction with
the lone pair of another small electronegative atom. Hydrogen-bond
energies are $0.1$ to $0.5\,\text{eV}$, intermediate between van der
Waals and covalent. Hydrogen bonds account for the anomalously high
boiling point of water (compared to H$_2$S, $-60\,\degC$), the
secondary structure of proteins, the double-helix structure of DNA,
and the unusual elastic and acoustic properties of ice and liquid
water.

\engterm{Dipole-dipole bonds}. Permanent dipoles in polar molecules
attract each other end-to-end. The interaction energy is comparable to
van der Waals and somewhat weaker than hydrogen bonding. It governs
the structure of polar liquids and the solubility patterns developed
in Chapter~3 of this volume.

Secondary bonds are weak compared to primary (covalent, ionic,
metallic) bonds, but they are central to the engineering of polymers
(Chapter~6 of this volume), of adhesives, of liquids, and of soft
biological tissues. Many engineering materials are mixtures: the
in-plane bonds of graphite are covalent, the inter-plane bonds are van
der Waals; the chains of polyethylene are covalent, the chains stick
to each other by van der Waals; the molecules of water are covalent,
the liquid is held together by hydrogen bonds. The reader who can
identify which class of bond is doing the work for which property
predicts the property correctly.

\begin{principle}[Bond character predicts property class]
The four canonical bond classes (ionic, covalent, metallic, secondary)
predict the broad property class of a solid. Ionic solids are hard,
brittle, electrically insulating in the solid, often transparent or
coloured by $d$-band transitions; covalent network solids are very
hard, electrically insulating or semiconducting, often optically clear
in the infrared; metals are ductile, electrically and thermally
conductive, optically reflective; molecular and polymeric solids
(secondary-bonded) are soft, low melting, often electrically
insulating. A material that does not fit one class clearly is a hybrid
or composite of two or more classes, and its properties are read by
identifying which bond does which job.
\end{principle}

\section{Crystals: lattices, defects, grains}\pathtag{standard}

A crystal is a periodic arrangement of atoms. The arrangement is
specified by a \engterm{Bravais lattice} (the set of translation
vectors that map the structure into itself) and a \engterm{basis} (the
atom or group of atoms attached to each lattice point). The fourteen
three-dimensional Bravais lattices, due to Auguste Bravais in 1848,
exhaust the possible crystalline geometries; engineering materials
draw on a small subset of them.

\engterm{Cubic lattices}. The simple cubic lattice (one atom per unit
cell at each corner of a cube) is rare in pure elements (polonium is
the only one). The body-centred cubic (BCC) lattice (one atom at each
corner plus one at the centre of the cell) is the room-temperature
structure of $\alpha$-iron, chromium, tungsten, molybdenum, vanadium,
sodium, potassium. The face-centred cubic (FCC) lattice (one atom at
each corner plus one at the centre of each face) is the structure of
aluminium, copper, gold, silver, nickel, platinum, lead, $\gamma$-iron
(above $912\,\degC$). The FCC lattice is the densest possible packing
of equal spheres (along with hexagonal close packing, HCP, which
shares the maximum packing fraction $\pi/\sqrt{18} \approx 0.74$).

\engterm{Hexagonal close-packed} (HCP). Two interpenetrating hexagonal
lattices, with the second offset to sit in the holes of the first.
The structure of zinc, magnesium, titanium ($\alpha$-phase below
$882\,\degC$), cobalt ($\alpha$-phase). The HCP packing fraction is
identical to FCC, but the stacking sequence is ABAB rather than ABCABC;
the difference matters for the dislocation behaviour developed in
Chapter~5 of this volume.

\engterm{Other lattices}. The diamond lattice (carbon, silicon,
germanium) is two interpenetrating FCC lattices, with each atom
tetrahedrally coordinated to four neighbours. The rock-salt structure
(NaCl, MgO, CaO, LiF) is two interpenetrating FCC lattices, with the
second shifted by half the cube edge. The zinc-blende structure (ZnS,
GaAs, InP) is the diamond structure with the two sublattices occupied
by different species. The wurtzite structure (ZnO, GaN, AlN) is the
same two-species arrangement on an HCP backbone. The perovskite
structure (CaTiO$_3$, BaTiO$_3$, lead zirconate titanate) is a
cubic arrangement with three distinct ion positions; it is the
structure of most ferroelectric ceramics.

\engterm{Crystallographic notation}. A direction in a cubic lattice is
written $[uvw]$, with $u, v, w$ integers giving the components of the
direction vector in units of the lattice constant. A plane is written
$(hkl)$, with $h, k, l$ the Miller indices (the reciprocals of the
plane's intercepts on the three axes, cleared of fractions). The
$[100]$ direction is along the cube edge; the $(100)$ plane is
perpendicular to it. The $[111]$ direction is along the cube diagonal;
the $(111)$ plane is perpendicular to it. The $\{hkl\}$ notation
denotes the family of planes equivalent to $(hkl)$ under the
crystal's symmetries; $\langle uvw \rangle$ denotes the family of
directions. The reader who works with crystals needs the notation to
read crystallographic data sheets; the engineering use of the notation
is for identifying slip planes and slip directions in plastic
deformation (Chapter~5).

\engterm{Defects}. Real crystals contain defects. The classification
is by dimensionality:

\begin{itemize}[leftmargin=*]
\item \engterm{Point defects} (zero-dimensional): vacancies (a missing
  atom from a lattice site), interstitials (an extra atom in a
  non-lattice position), substitutional impurities (a foreign atom
  occupying a lattice site), interstitial impurities (a foreign atom
  in a non-lattice position). The equilibrium concentration of
  vacancies at temperature $T$ is $n_v/n = \exp(-Q_v/k_B T)$, with
  $Q_v$ the vacancy formation energy ($\sim 1\,\text{eV}$ for typical
  metals). For copper at room temperature, $n_v/n \sim 10^{-17}$; at
  $1000\,\degC$, $\sim 10^{-4}$. Vacancies enable diffusion
  (Chapter~4 of this volume) and contribute to high-temperature creep
  (Chapter~5).
\item \engterm{Line defects} (one-dimensional): dislocations, the
  carriers of plastic deformation in crystalline solids. The edge
  dislocation is an extra half-plane of atoms inserted into the
  crystal; the screw dislocation is a helical winding of lattice
  planes around the dislocation line. Dislocations are characterised
  by their Burgers vector $\vec{b}$ (the closure failure of an atomic
  circuit around the dislocation core), and they move under shear
  stress, producing plastic strain at stresses far below the
  theoretical shear strength of a perfect crystal. The dislocation
  theory of plastic deformation, developed in the 1930s by Taylor,
  Orowan, and Polanyi, is the bridge between atomic bonding and
  macroscopic plasticity \cite[ch.~6]{text:smallman-physical-metallurgy}.
  Chapter~5 of this volume develops the theory and its consequences
  for strengthening mechanisms.
\item \engterm{Surface defects} (two-dimensional): grain boundaries
  (the surfaces between crystals of different orientation in a
  polycrystalline solid), twin boundaries (surfaces across which the
  lattice is mirror-imaged), stacking faults (a local error in the
  stacking sequence of close-packed planes). Grain boundaries dominate
  the room-temperature mechanical behaviour of most engineering
  metals: a fine-grained material is stronger than a coarse-grained
  material of the same composition, according to the Hall-Petch
  relation $\sigma_y = \sigma_0 + k d^{-1/2}$ with $d$ the average
  grain diameter (Chapter~5).
\item \engterm{Volume defects} (three-dimensional): voids, inclusions,
  precipitates, second-phase particles. Precipitates are the central
  strengthening mechanism in age-hardened aluminium and steel
  (Chapter~5); inclusions are common origins of fatigue crack
  initiation (Chapter~9).
\end{itemize}

\engterm{Grains}. A polycrystalline metal is an aggregate of many
small single crystals (\engterm{grains}), each with its own
orientation, separated by grain boundaries. Grain sizes in typical
engineering metals range from $1\,\si{\micro\meter}$ in
heavily-cold-worked or rapidly-solidified material to $1\,
\text{mm}$ in slowly-cooled cast structures. The grain structure is
revealed by polishing and etching a metal sample and viewing it under
an optical microscope; the etchant attacks the grain boundaries faster
than the grain interiors, leaving a contrast network. The grain
structure is one of the most easily measured features of a metal and
one of the most predictive of its mechanical behaviour.

\section{Amorphous solids}\pathtag{standard}

An amorphous solid lacks long-range periodic order. Short-range order
(the nearest-neighbour coordination, set by the bond geometry) is
preserved; medium-range order (the second and third neighbour shells)
is preserved approximately; long-range order is absent. The classic
amorphous materials are silica glass (SiO$_4$ tetrahedra connected in
a random network), the organic polymers in their non-crystalline
state, and the metallic glasses produced by quenching certain alloy
compositions faster than they can crystallise.

\engterm{Glass formation}. A liquid cooled below its melting point can
take one of two routes. If it nucleates and grows crystals, it becomes
a crystalline solid; if it does not, it becomes a supercooled liquid
and, on further cooling, a glass. The transition from supercooled
liquid to glass is the \engterm{glass transition}, a temperature
$T_g$ at which the relaxation time of the liquid becomes too long for
the liquid to equilibrate on laboratory timescales. The glass
transition is not a thermodynamic transition in the strict sense (no
discontinuity in the entropy or volume), but a kinetic one: it
depends on the cooling rate, with faster cooling producing a higher
$T_g$. The shift is typically a few kelvins per decade of cooling-rate
variation. For window glass, $T_g \approx 1000\,\si{\kelvin}$; for
polystyrene, $\approx 370\,\si{\kelvin}$; for natural rubber, $\approx
210\,\si{\kelvin}$. Chapter~6 of this volume develops the glass
transition for polymers.

\engterm{Properties of amorphous solids}. Without long-range order,
amorphous solids have no cleavage planes (they fracture along
conchoidal surfaces in glass, in random patterns in metallic glass),
no Bragg diffraction peaks (X-ray diffraction shows broad halos
rather than sharp lines), and isotropic macroscopic properties (no
preferred direction). The mechanical behaviour of amorphous solids
above $T_g$ is liquid-like and below $T_g$ is glassy; the transition
is smooth on the macroscopic timescale and abrupt on the molecular.

\engterm{Engineering amorphous materials}. Silica glasses and
soda-lime glasses are the dominant transparent structural materials;
polymer glasses (polystyrene, polycarbonate, PMMA) cover the
lightweight transparent applications; metallic glasses (Pd-Si, Zr-Cu,
Fe-B based compositions) have unusual combinations of strength and
elastic resilience because they have no grain boundaries or
dislocations to mediate plastic deformation, but they are mechanically
brittle for the same reason. The amorphous semiconductor amorphous
silicon (a-Si:H, with hydrogen passivation of dangling bonds) is the
working material of low-cost thin-film photovoltaics and of liquid-
crystal-display backplanes. The reader who specifies an amorphous
material chooses it for properties (transparency, isotropy, absence
of grain boundaries) that the crystalline form does not provide.

\section{Why bonding determines properties}\pathtag{core}

The bonding budget of a material predicts its principal engineering
properties to within a constant factor. The connection is the topic
that earlier sections of this chapter have been building toward; the
present section makes it explicit.

\engterm{Elastic stiffness} ($E$, Young's modulus). The energy per
atom of a bond at its equilibrium length is approximately $U_0$ (a
typical bond energy, $1$ to $10\,\text{eV}$). The bond's stiffness,
the force per unit displacement near equilibrium, is approximately
$k \sim U_0/r_0^2$ with $r_0$ the equilibrium bond length. The
macroscopic Young's modulus follows by dividing the force per atom by
the area per atom:
\[
E \sim \frac{k}{r_0} \sim \frac{U_0}{r_0^3}.
\]
For $U_0 \sim 3\,\text{eV}$ and $r_0 \sim 0.2\,\text{nm}$ the estimate
is $E \sim 600\,\text{GPa}$, comparable to the measured Young's
modulus of tungsten ($410\,\text{GPa}$), diamond ($1100\,\text{GPa}$),
and the upper end of the engineering ceramics. For weaker secondary
bonds ($U_0 \sim 0.1\,\text{eV}$ at $r_0 \sim 0.3\,\text{nm}$),
$E \sim 6\,\text{GPa}$, comparable to the bulk modulus of glassy
polymers. Young's modulus thus ranges over three orders of magnitude
across material classes; the variation tracks the bond strength
directly. The Ashby chart of $E$ versus density (Chapter~10 of this
volume) makes the patterns visible at a glance.

\engterm{Theoretical strength}. The stress to break a bond is
approximately $\sigma_{\text{theor}} \sim E/10$ to $E/5$. The
factor of order $10$ comes from the curvature of the interatomic
potential and is set by the bond's anharmonicity. For diamond
($E = 1100\,\text{GPa}$), $\sigma_{\text{theor}} \sim 100\,\text{GPa}$;
for polyethylene ($E \sim 1\,\text{GPa}$), $\sigma_{\text{theor}} \sim
100\,\text{MPa}$. The measured strengths are lower than the
theoretical, by factors of $10$ to $1000$, because real materials
contain defects that initiate failure before the bond limit is
reached. The gap between theoretical and observed strength is the
subject of fracture mechanics (Volume III, Ch~7) and of the
strengthening-mechanism literature in metallurgy (Chapter~5).

\engterm{Melting point}. The energy required to disrupt the bonded
arrangement of a solid into the disordered arrangement of a liquid
scales with the bond energy. The empirical rule that $T_m \sim
U_0/k_B$ (with a numerical factor of $0.05$ to $0.5$ across material
classes) recovers melting points within a factor of three for most
materials. Diamond ($U_0 = 3.6\,\text{eV}$, $T_m \approx 4300\,
\si{\kelvin}$); silicon ($U_0 = 2.3\,\text{eV}$, $T_m = 1683\,
\si{\kelvin}$); copper ($U_0 = 3.5\,\text{eV}$ cohesive, $T_m = 1358\,
\si{\kelvin}$); polyethylene secondary bonds dominate ($U_0 \sim 0.1\,
\text{eV}$ per CH$_2$, $T_m \approx 410\,\si{\kelvin}$).

\engterm{Thermal expansion}. The asymmetry of the interatomic
potential (anharmonicity) gives the bond a longer mean length at
higher temperature. The linear thermal expansion coefficient $\alpha$
is approximately
\[
\alpha \sim \frac{k_B}{U_0 r_0},
\]
which gives $\alpha \sim 10^{-5}\,/\si{\kelvin}$ for typical
$U_0 \sim 3\,\text{eV}$ and $r_0 \sim 0.2\,\text{nm}$. The strongly-
bonded ceramics (diamond, silicon carbide) have $\alpha \sim 3 \times
10^{-6}\,/\si{\kelvin}$; metals $\sim 10^{-5}\,/\si{\kelvin}$;
polymers $\sim 10^{-4}\,/\si{\kelvin}$. The thermal-expansion ranking
tracks the bond-energy ranking.

\engterm{Electrical conductivity}. Determined by the electronic
structure, not by the bond geometry alone. Metals (delocalised
electrons) conduct; semiconductors (small gap between filled valence
band and empty conduction band) conduct weakly and strongly
temperature-dependent; insulators (large gap) do not conduct. The
band structure is set by the bond type and the crystal symmetry; for
the engineering reader, the rule of thumb is that strong covalent or
ionic bonding produces large band gaps and insulating behaviour
unless the material is doped, while metallic bonding produces
conductors. Volume VII Chapter~3 develops the band theory.

\engterm{Thermal conductivity}. In metals, dominated by the same
electron sea that carries electrical current. In insulators, dominated
by phonons (quantised lattice vibrations); the thermal conductivity
is set by the speed of sound and the phonon mean free path, with
defects and grain boundaries acting as phonon scatterers. Diamond is
an outlier: $2000\,\text{W}/(\text{m}\cdot\si{\kelvin})$ at room
temperature, the highest of any solid, because the stiff covalent
lattice supports high-frequency phonons with long mean free paths
despite being an electrical insulator.

The chapter's organising claim is that the engineer can predict a
material's properties to within a factor of three by reading off its
position in the periodic table, identifying its dominant bond type,
and applying the scaling relations above. Volume V refines the
estimates by adding processing and microstructure, but the bond-level
scaffolding survives every later refinement.

\begin{estimation}
\textbf{Question.} Estimate Young's modulus of a single crystal of
sodium chloride at room temperature.

\smallskip
\textit{Estimate, before reading on.}

\smallskip
The bond is ionic, between Na$^+$ and Cl$^-$ ions. The lattice energy
per ion pair is approximately $-7.9\,\text{eV}$ from the Madelung
calculation above. The nearest-neighbour distance is $r_0 = 2.82\,
\si{\angstrom} = 2.82 \times 10^{-10}\,\text{m}$.

The bond stiffness per pair is
\[
k \sim \frac{|U_0|}{r_0^2} = \frac{7.9 \times 1.6 \times 10^{-19}\,\text{J}}{(2.82
\times 10^{-10})^2\,\text{m}^2} \approx 16\,\text{N}/\text{m}.
\]

Young's modulus is the bond stiffness divided by the area per atom
times the bond length:
\[
E \sim \frac{k}{r_0} = \frac{16}{2.82 \times 10^{-10}} \approx 5.7 \times
10^{10}\,\text{Pa} = 57\,\text{GPa}.
\]

The measured Young's modulus of single-crystal NaCl along the
$\langle 100 \rangle$ direction is $44\,\text{GPa}$, within $30\,\%$
of our estimate.

\smallskip
\textit{Postmortem.} The estimate was honest about the bond energy
(we used the lattice energy without correcting for the partial
contribution of nearest-neighbour bonds), about the bond length, and
about the geometric prefactor. The agreement with the measured value
to within $30\,\%$ is the expected accuracy of the bond-level scaling
argument. The reader who applies the same calculation to MgO (lattice
energy $-3900\,\text{kJ}/\text{mol}$, $r_0 = 2.10\,\si{\angstrom}$,
estimated $E \sim 380\,\text{GPa}$, measured $E = 300\,\text{GPa}$)
gets the same accuracy. The bond-budget estimate of Young's modulus
is one of the easiest checks an engineer can perform on the
plausibility of a manufacturer's data sheet.
\end{estimation}

\section{Failure: bond-level origins of brittleness, ductility, conductivity}\pathtag{core}

The macroscopic failure modes of a material trace to its bond-level
behaviour at the point of failure. The chapter's failure section names
three such linkages.

\engterm{Brittle fracture in ionic and covalent solids}. An ionic
lattice in shear translates one plane of ions over the next; at
intermediate displacements like-charged ions face each other and the
electrostatic repulsion exceeds the attraction. The lattice cleaves
along the plane of like-ion approach rather than sliding plastically.
The macroscopic phenomenon is brittle fracture along low-index
crystallographic planes, the visible cleavage faces on a fractured
ionic single crystal. Covalent network solids fracture by the same
brittle mechanism but for a different bond-level reason: the
directional covalent bonds resist shear translation because the
deflected bond angles incur a large energy cost. Engineering
consequences are visible across the volume: ceramic teacups crack
rather than dent (Chapter~7); silicon wafers split along $\{111\}$
cleavage planes (Volume VII Ch~5); the windshield of a car spreads a
visible crack network because the laminated safety glass is
ionic-bonded soda-lime silica with a polymer interlayer (Chapter~7
and Chapter~8).

\engterm{Ductility in metals}. The metallic bond, being non-directional
and tolerant of geometric rearrangement, allows the lattice to slip
under shear without breaking specific localised bonds. The
microscopic carrier of the slip is the dislocation; the macroscopic
phenomenon is plastic deformation at modest stress, with strain to
failure of $10\,\%$ to $50\,\%$ or more for ductile metals. The
dislocation theory of plastic deformation, developed in detail in
Chapter~5, is the bridge between the metallic bond's non-directionality
and the engineering ductility of structural steels and aluminium
alloys. The engineering consequence is that metals deform visibly
before they fail, giving an inspecting engineer a warning sign that
ceramic and glass components do not provide.

\engterm{Conductivity and its loss}. Metallic conductivity is the
free flow of conduction electrons under an applied field. The
conductivity is reduced by anything that scatters electrons:
thermal vibrations (phonon scattering, giving the $T^1$ resistivity of
pure metals at high temperature), impurity atoms (residual resistivity
at low temperature), grain boundaries (resistivity in
nanocrystalline materials), and dislocations and other defects. The
engineering consequence is that a metal's conductivity is a sensitive
probe of its purity and microstructure: a copper wire whose
resistivity exceeds the handbook value by $5\,\%$ contains alloying
elements or work-hardening defects in amounts the engineer can
estimate. Conductivity loss in service is a failure mode in its own
right: oxide films on contacts (Chapter~9), grain growth in
high-temperature service, and radiation damage in nuclear environments
all increase resistivity and degrade the electrical function.

The failure modes recur in later chapters of the volume in
mechanism-specific form. The present section's claim is that the
fundamental story is set at the bond level: brittle solids fail by
bond rupture against shear-induced repulsion, ductile solids fail by
dislocation propagation through a bond-tolerant lattice, and
conductors lose conductivity when scattering centres interrupt the
electron sea.

\begin{principle}[Property failure mirrors bond character]
A material's principal failure mode mirrors its dominant bond
character. Ionic and covalent solids fail brittly by bond rupture;
metallic solids fail ductilely by dislocation motion until
work-hardening or void coalescence exhausts their plastic capacity;
secondary-bonded solids creep, soften, and disentangle under load
because the secondary bonds can rearrange thermally. The engineer
predicting the failure mode of a new material identifies the bond
mix and applies the appropriate failure model from the relevant
chapter of this volume.
\end{principle}

\section{Project}\pathtag{core}

\begin{project}[Five household materials, three properties each]
The reader identifies five household materials by their dominant
bonding type and explains three of the materials' macroscopic
properties accordingly.

\textbf{Track}. Household identification and explanation.
\textbf{Hazard class}: low. No demolition, no chemical reagents
beyond the materials themselves.

\textbf{Selection rule}. The five materials must span at least four of
the four bond classes (ionic, covalent network, metallic, secondary).
Candidate household materials:

\begin{itemize}[leftmargin=*]
\item Ceramic mug or unglazed terracotta flower pot (ionic, possibly
  with covalent network components).
\item Window glass or eyeglass lens (amorphous covalent network).
\item Stainless-steel cutlery or copper pipe (metallic).
\item Polyethylene shopping bag or polypropylene food container
  (long covalent chains, secondary bonds between chains).
\item Wax candle or block of soap (van der Waals plus weak
  hydrogen bonds).
\end{itemize}

\textbf{Procedure}.

\begin{enumerate}[leftmargin=*]
\item For each of the five materials, write a one-paragraph entry that
  identifies its dominant bond type, names the elemental composition
  (approximate is acceptable), and locates the constituent elements on
  the periodic table.
\item For each material, identify three macroscopic properties that
  the reader can observe or measure with household tools: density (by
  weighing and water displacement); a stiffness comparison (which of
  the five is hardest to scratch with a fingernail or a coin); a
  thermal-response comparison (which warms fastest to the touch when
  held); transparency (yes or no); electrical conductivity (touch
  with a multimeter probe).
\item Explain each measured property in terms of the bond type,
  using the scaling relations from Section~6.
\item Identify one property prediction that the bond-level argument
  gets wrong, and propose a refinement (composition, microstructure,
  defect population) that explains the discrepancy.
\end{enumerate}

\textbf{Deliverable}. A five- to ten-page report with one table
summarising the five materials and one bond-property comparison
chart. Include a paragraph identifying the prediction that failed and
the proposed refinement.

\textbf{Time}. Two to four hours of observation, two to four hours of
write-up.
\end{project}

\section{Exercises}

\subsection*{Calculation}\pathtag{core}

\begin{exercise}
Compute the energy in joules of one mole of C-C bonds (bond energy
$3.6\,\text{eV}$ per bond). Compare to the energy of one mole of
hydrogen-bond pairs ($0.2\,\text{eV}$ per bond) and to one mole of
van der Waals pairs ($0.1\,\text{eV}$).
\end{exercise}

\begin{exercise}
The lattice energy of NaCl is approximately $-787\,\text{kJ}/\text{mol}$.
Estimate the Coulomb contribution to this energy from
$U_{\text{Coul}} = -M e^2 N_A/(4\pi\varepsilon_0 r_0)$ with the Madelung
constant $M = 1.7476$ and $r_0 = 2.82\,\si{\angstrom}$. Compare to the
measured value and account for the Born repulsion term using the
factor $(1 - 1/n)$ with $n = 8$.
\end{exercise}

\begin{exercise}
Determine the linear density of atoms (atoms per metre) along the
$[100]$, $[110]$, and $[111]$ directions of an FCC copper crystal
with lattice constant $a = 3.61\,\si{\angstrom}$.
\end{exercise}

\begin{exercise}
Compute the planar density of atoms (atoms per square metre) on the
$(100)$, $(110)$, and $(111)$ planes of a BCC iron crystal with
lattice constant $a = 2.87\,\si{\angstrom}$. Identify which plane has
the highest planar density.
\end{exercise}

\begin{exercise}
The vacancy formation energy in copper is $Q_v = 0.9\,\text{eV}$.
Compute the equilibrium fraction of vacant sites at room temperature
($300\,\si{\kelvin}$) and at $1000\,\degC$. Use $k_B = 8.617 \times
10^{-5}\,\text{eV}/\si{\kelvin}$.
\end{exercise}

\begin{exercise}
The Fermi energy of aluminium is $E_F = 11.7\,\text{eV}$ at zero
temperature. Aluminium contributes three valence electrons per atom.
Verify the Fermi energy from the free-electron formula $E_F =
(\hbar^2/2 m_e)(3\pi^2 n)^{2/3}$ given the density $\rho = 2700\,
\text{kg}/\text{m}^3$ and the atomic mass $M_a = 27.0\,\text{g}/\text{mol}$.
\end{exercise}

\begin{exercise}
Estimate the Young's modulus of diamond from $E \sim U_0/r_0^3$, with
$U_0 = 3.6\,\text{eV}$ and $r_0 = 1.54\,\si{\angstrom}$. Compare to
the measured value of approximately $1100\,\text{GPa}$.
\end{exercise}

\begin{exercise}
Compute the packing fraction of the simple-cubic, body-centred-cubic,
and face-centred-cubic lattices, assuming rigid spheres in contact
along the close-packed direction.
\end{exercise}

\subsection*{Derivation}\pathtag{standard}

\begin{exercise}
Starting from the Bohr atomic model, derive the formula $E_n = -13.6\,
\text{eV}/n^2$ for the energy of the hydrogen atom in the $n$-th
quantum state. State the assumptions and identify which fail for
multi-electron atoms.
\end{exercise}

\begin{exercise}
Derive the Madelung-sum form of the lattice energy of an ionic crystal
$U = -M e^2/(4\pi\varepsilon_0 r_0)$ for the case of NaCl, by
summing the Coulomb pair energies over successive shells of neighbours.
Discuss the convergence of the series.
\end{exercise}

\begin{exercise}
Derive the scaling $E \sim U_0/r_0^3$ for the Young's modulus of a
bonded solid by linearising the interatomic potential about its
equilibrium and dividing the bond stiffness by the area per atom.
\end{exercise}

\begin{exercise}
Derive the Hall-Petch relation $\sigma_y = \sigma_0 + k d^{-1/2}$
qualitatively from a dislocation-pileup model in which the
stress concentration at a grain boundary scales as the square root of
the pile-up length.
\end{exercise}

\subsection*{Estimation}\pathtag{core}

\begin{exercise}
Estimate the number of atoms in a one-centimetre cube of copper.
\end{exercise}

\begin{exercise}
Estimate the thermal-expansion coefficient of a polymer ($U_0 \sim
0.1\,\text{eV}$, $r_0 \sim 0.3\,\text{nm}$) and of a ceramic ($U_0
\sim 5\,\text{eV}$, $r_0 \sim 0.2\,\text{nm}$). Compare the ratio to
the measured values for polyethylene and alumina.
\end{exercise}

\begin{exercise}
Estimate the surface energy of a NaCl crystal (energy per unit area
of a freshly cleaved $\{100\}$ face) from the ionic bond energy and
the planar density of broken bonds at the cleavage plane.
\end{exercise}

\subsection*{Simulation}\pathtag{mastery}

\begin{exercise}
Implement a one-dimensional chain of atoms connected by Lennard-Jones
springs ($V(r) = 4\varepsilon[(\sigma/r)^{12} - (\sigma/r)^6]$).
Compute the equilibrium spacing, the spring constant near
equilibrium, and the anharmonic shift of the mean spacing with
temperature.
\end{exercise}

\begin{exercise}
Simulate a two-dimensional FCC crystal with a single edge dislocation.
Apply a uniform shear stress and observe the motion of the
dislocation. Compute the Peierls stress (the minimum stress for
dislocation motion).
\end{exercise}

\subsection*{Design}\pathtag{standard}

\begin{exercise}
Specify a candidate material class for a structural component that
must combine high stiffness ($E > 200\,\text{GPa}$), low density
($\rho < 5000\,\text{kg}/\text{m}^3$), and ductile fracture (strain
to failure $> 0.05$). Identify which bond type dominates and which
elements are candidates.
\end{exercise}

\begin{exercise}
Specify a transparent material for a flat-panel display cover that
must be hard ($H > 5\,\text{GPa}$ Vickers) and resistant to scratching
by sand grains. Identify the bond type and propose one candidate
chemistry and one geometric strategy (multilayer, surface treatment)
for the design.
\end{exercise}

\begin{exercise}
Specify the electrical-conductor material for a long-distance
overhead power transmission line. State the bond-level argument for
the choice and the practical considerations (price, ductility,
oxidation) that the bond-level argument does not capture.
\end{exercise}

\subsection*{Diagnosis and reverse engineering}\pathtag{core}

\begin{exercise}
A consumer ceramic mug breaks into a small number of sharp pieces
when dropped on a tile floor; a polystyrene cup of the same size,
dropped from the same height, deforms slightly and bounces. Identify
the bond type of each material and explain the difference in
fracture mode in terms of the dislocation availability and the bond
directionality.
\end{exercise}

\begin{exercise}
A copper wire that should carry $20\,\text{A}$ heats up noticeably
along part of its length and stays cool elsewhere. The wire is uniform
in diameter. Identify the most likely bond-level cause of the
non-uniform resistivity and propose one diagnostic measurement.
\end{exercise}

\begin{exercise}
A piece of window glass and a piece of metallic glass (both amorphous)
look superficially similar in transparency. Identify two macroscopic
tests that would distinguish the two unambiguously.
\end{exercise}

\subsection*{Failure analysis}\pathtag{core}

\begin{exercise}
A single-crystal silicon wafer cracks along $\{111\}$ planes when
notched. Identify the bond geometry that produces the cleavage and
estimate the theoretical cleavage stress from $\sigma \sim E/10$.
\end{exercise}

\begin{exercise}
A polymer adhesive bond between two metal plates fails by clean
separation at the polymer-metal interface, with no visible polymer
on the metal. Identify which class of bond was acting at the
interface and propose a surface treatment that would convert the
failure to a cohesive failure within the polymer.
\end{exercise}

\subsection*{Judgment}\pathtag{standard}

\begin{exercise}
A colleague proposes that a new alloy must be metallic because it
conducts electricity. The alloy is in fact a heavily-doped
semiconductor. Identify the additional measurements required to
distinguish a true metal from a degenerately doped semiconductor by
the temperature dependence of conductivity.
\end{exercise}

\begin{exercise}
A vendor advertises ``ionic-bonded steel'' as a structural product.
Write a one-paragraph reply identifying the category error and the
correct bond description of steel.
\end{exercise}

\begin{exercise}
A reader argues that the periodic table is ``just an organising
chart'' and adds no information beyond what the elemental list
already provides. Write a one-paragraph reply naming two predictions
that the table makes about element behaviour and that an
alphabetised list cannot make.
\end{exercise}
